\section{Problem Summary}
You should receive every integer from \(1\) through \(n\), except that one value is missing. The task is to determine which value never appears.

\section{Editorial}
The full set \(1 + 2 + \cdots + n\) has a known sum:
\[
\frac{n(n + 1)}{2}.
\]

If we subtract the sum of the numbers that were actually provided, the remainder must be the missing value.

That avoids sorting, marking arrays, or any more complicated bookkeeping. A single pass through the input is enough.

\section{Complexity Analysis}
\begin{itemize}
\item Time: \(O(n)\)
\item Memory: \(O(1)\)
\end{itemize}

\section{Notes}
Compute the sums with 64-bit integers. The formula can overflow a 32-bit type even though the answer itself is small.
